% cSpell:ignore Schmitt,monostable,multivibrator,555,7413,74121,CMOS,TTL,VCC,LF398,ADC,DAC,CB7520,DAC0808,AD9012,AD9002,AD9020,ADC0804
\setlength{\parskip}{0.4em}
\setlength{\itemsep}{0.2em}
\setlength{\tabcolsep}{4pt}

\chapter{数一模和模一数转换}
\section{概述}
\subsection{应用用途}
计算机控制系统中，被控对象多为温度、压力、位移等非电量；传感器将非电量转为连续模拟量，经A/D转换得到数字量送入计算机运算处理；计算机输出数字量再经D/A转换还原模拟量，驱动执行机构调节被控对象。

\subsection{核心性能指标}
转换精度、转换速度是DAC与ADC两类转换器最关键性能指标。

\subsection{基本概念与分类}
1. \textbf{D/A转换器 DAC}
数字量向模拟量转换电路；主流类型：权电阻网络DAC、倒T形电阻网络DAC、权电流型DAC、权电容型DAC、开关树型DAC。
2. \textbf{A/D转换器 ADC}
模拟量向数字量转换电路；分为直接型ADC、间接型ADC；按输出格式分为并行输出ADC、串行输出ADC。
教材先讲解DAC，DAC电路原理更简单，也是ADC内部组成单元。

\section{D/A转换器}
\subsection{DAC整体结构框架}
完整n位DAC由五部分构成：\textbf{数码寄存器、模拟开关阵列、解码电阻网络、求和放大电路、基准电压源}。
数字量串行/并行存入寄存器，寄存器输出逐位驱动模拟开关，开关将对应权值电流送入求和放大器，叠加得到与输入数字量成正比的模拟输出。

\subsection{DAC分类体系}
1. 按解码网络结构：T型电阻DAC、倒T型电阻DAC、权电流DAC、权电阻DAC；
2. 按模拟开关类型：
    - CMOS开关型（功耗低，转换速度慢）；
    - 双极型电流开关型（转换速度快，包含ECL电流开关）。

\subsection{权电阻网络D/A转换器}
\subsubsection{二进制数字量权值表达}
n位二进制数 $D_n = d_{n-1}d_{n-2}\dots d_1d_0 = 2^{n-1}d_{n-1}+2^{n-2}d_{n-2}+\dots+2^1d_1+2^0d_0$
$2^{n-1}$ 为最高位MSB权值，$2^0$ 为最低位LSB权值。

\subsubsection{电路组成}
4位权电阻DAC由：权电阻网络、四路模拟开关、求和运算放大器、基准电压$V_\text{REF}$组成。
开关逻辑：$d_i=1$ 开关接$V_\text{REF}$产生支路电流；$d_i=0$ 开关接地，支路电流为0。
支路权电流（4位）：
\[
I_3=\frac{V_\text{REF}}{R},\; I_2=\frac{V_\text{REF}}{2R},\; I_1=\frac{V_\text{REF}}{4R},\; I_0=\frac{V_\text{REF}}{8R}
\]

\subsubsection{运放虚短虚断基础}
深度负反馈运放线性区满足：虚短 $V_+=V_-=0$；虚断输入电流为0；反相比例输出 $v_O = -\frac{R_f}{R_1}U_1$。

\subsubsection{输出电压推导（取$R_f=R/2$）}
总求和电流 $i_\Sigma = I_3+I_2+I_1+I_0$
\[
\begin{aligned}
v_O &= -R_f i_\Sigma
= -\frac{R}{2}\left( \frac{V_\text{REF}}{R}d_3+\frac{V_\text{REF}}{2R}d_2+\frac{V_\text{REF}}{4R}d_1+\frac{V_\text{REF}}{8R}d_0 \right) \\
&= -\frac{V_\text{REF}}{2^4}\left(2^3d_3+2^2d_2+2^1d_1+2^0d_0\right)
\end{aligned}
\]
通用n位公式：
\[
v_O = -\frac{V_\text{REF}}{2^n} D_n
\]
输入全0时$v_O=0$；输入全1时 $v_O=-\frac{2^n-1}{2^n}V_\text{REF}$，输出范围 $0 \sim -\frac{2^n-1}{2^n}V_\text{REF}$。

\subsubsection{优缺点}
优点：电路结构简单、元件数量少；
缺点：各位电阻阻值相差极大，位数越多阻值差距越大，高精度电阻生产困难。

\subsection{倒T形电阻网络D/A转换器}
\subsubsection{设计目的}
解决权电阻DAC电阻值差异过大、精度难以保证的缺陷，仅使用$R、2R$两种电阻。

\subsubsection{电路电流特性}
利用虚地 $V_- \approx V_+=0$，无论开关接运放输入端还是地，每条2R支路流过电流恒定，从MSB到LSB电流逐次二分。
总基准电流 $I=\frac{V_\text{REF}}{R}$。

\subsubsection{n位通用输出公式}
\[
v_O = -\frac{V_\text{REF}}{2^n}\left(2^{n-1}d_{n-1}+\dots+2^1d_1+2^0d_0\right)
= -\frac{V_\text{REF}}{2^n}D_n
\]
输出范围、比例关系与权电阻DAC完全一致；$V_\text{REF}$取负值可得到正向输出电压。

\subsubsection{集成DAC CB7520（AD7520）}
10位倒T型CMOS开关DAC，片内集成倒T型$R/2R$电阻网络，需要外接运算放大器；基准电压$V_\text{REF}$精度直接决定转换精度。
10位输出公式：
\[
v_O = -\frac{V_\text{REF}}{2^{10}} D_{10}
\]

\subsubsection{单极性输出外围电路}
可外接调零电阻、反馈调节电阻，调节运放增益改变满量程输出幅度；$R_{W3}$为运放调零电阻，$R_{W1}$调节放大比例，$R_{W2}$串联分压缩小输入基准电流。

\subsection{权电流型D/A转换器}
\subsubsection{改进思路}
权电阻、倒T型DAC开关导通压降会引入转换误差；权电流DAC采用恒流源支路，支路电流不受开关导通内阻压降影响，降低开关匹配要求。

\subsubsection{电路原理}
恒流源支路电流仅由发射极回路参数固定，数字位$d_i=1$恒流源接入运放求和端，$d_i=0$接地；搭配倒T型电阻网络减少电阻种类。

\subsubsection{典型集成芯片 DAC0808}
8位双极型电流开关DAC，转换速度高；内置恒流基准电路，外部仅需运放、补偿电容、基准电压。

\subsection{双极性输出D/A转换器}
数字量包含正负符号（补码），需要DAC输出正负模拟电压；在单极性倒T型DAC求和端叠加偏移基准电流，偏移电流大小等于最高位权电流，实现正负对称输出。

\subsection{DAC转换精度与转换速度}
\subsubsection{转换精度两大指标}
1. \textbf{分辨率（理论精度）}
表征DAC分辨输入微小数字变化的能力，总共有$2^n$个输出等级：
\[
\text{分辨率} = \frac{1}{2^n - 1}
\]
位数$n$越大分辨率越高，也可用可分辨最小电压与满量程比值表示。

2. \textbf{转换误差（实际精度）}
实际输出与理想理论值的最大偏差，由基准源漂移、运放零点、开关内阻、电阻失配、器件非线性共同导致；分为三类：
    a) 比例系数误差：基准电压$V_\text{REF}$偏移带来，误差与输入数字成正比；
    b) 失调误差（平移误差）：运放零点漂移，整条转换曲线整体上下平移，与输入数字无关；
    c) 非线性误差：开关导通电阻、电阻失配引起，误差无固定线性规律；
总误差为各类误差绝对值叠加。

\subsubsection{转换速度——建立时间$t_\text{set}$}
数字输入突变后，输出模拟量进入稳态误差$\pm\frac{1}{2}\text{LSB}$区间所需时长；数字跳变幅度越大建立时间越长。
不含运放DAC：$t_\text{set}<0.1\mu\text{s}$；集成运放DAC：$t_\text{set}>1.5\mu\text{s}$；高速系统选用高速运放缩短建立时间。

\section{A/D转换器}
\subsection{A/D转换完整流程与前置单元}
\subsubsection{四大核心步骤}
连续模拟量转换为离散数字量，必须依次完成\textbf{取样、保持、量化、编码}四个环节，工程中取样和保持常集成在同一电路实现。

\subsubsection{取样定理（奈奎斯特采样定理）}
若输入模拟信号最高频率分量为$f_{i(\text{max})}$，取样频率$f_s$必须满足：
\[
f_s \ge 2f_{i(\text{max})}
\]
工程实用取值：$f_s=(3\sim5)f_{i(\text{max})}$；每次采样后必须保持采样电平稳定一段时间，为后续量化编码提供稳定输入。

\subsubsection{取样—保持电路}
1. \textbf{基础电路结构}
由运算放大器$\text{A}$、保持电容$C_H$、NMOS模拟开关$\text{T}$构成，$v_I$为模拟输入，$v_L$为取样控制时钟信号。
- 取样阶段（$v_L$高电平）：MOS管导通，输入电压通过$R_1$、导通管$\text{T}$向$C_H$充电，若取$R_1=R_f$且运放理想，则$v_O=v_C=-v_I$；
- 保持阶段（$v_L$低电平）：MOS管截止，电容$C_H$无放电回路，两端电压维持不变，输出锁定采样时刻电平；
保持电容漏电流越小、运放输入阻抗越高，保持时长越长。

2. \textbf{集成取样保持芯片LF398}
内部包含两级运放$\text{A}_1/\text{A}_2$、逻辑控制单元、模拟开关与保护二极管；$v_\text{OS}$为失调调零端，可外接电阻修正输出零点；二极管$D_1/D_2$用于过压保护，避免开关单元承受过高电压。

\subsubsection{量化与编码}
1. \textbf{量化定义}
将连续模拟电压划分成多个离散等级，每一级最小电压间隔称为量化单位$\Delta$（对应数字量最低位LSB），所有采样电压归并到最近的量化等级，该过程会引入固有量化误差。
2. \textbf{两种量化方式}
- \textbf{只舍不入量化}：不足一个量化单位的部分全部舍去，最大量化误差为$\Delta$；
- \textbf{四舍五入量化}：小于半个量化单位舍去、大于等于半个量化单位向上取整，最大量化误差仅$\Delta/2$，工业ADC普遍采用该方案。
3. \textbf{编码}
将量化后的离散电平用二进制（或其他进制）代码表示，最终输出标准数字信号；正负范围输入多采用二进制补码编码。

\subsubsection{ADC整体分类框架}
\begin{itemize}
\item 直接型ADC：并联比较型、反馈比较型（计数型、逐次渐近型）；
\item 间接型ADC：V-T变换型（单积分、双积分、四积分）、V-F变换型；
\end{itemize}

\subsection{直接型A/D转换器}
\subsubsection{并联比较型ADC}
1. \textbf{电路组成}
由分压电阻链、电压比较器组、寄存器、代码转换逻辑三部分构成，属于高速并行直接ADC。
2. \textbf{量化规则}
满量程$0\sim V_\text{REF}$范围内，$n$位ADC量化单位：
\[
\Delta=\frac{2V_\text{REF}}{2^{n+1}}
\]
比较器根据输入电平输出高低电平，经寄存器锁存后由组合逻辑转换为标准二进制输出。
3. \textbf{核心优缺点}
优点：转换速度极快，单次转换仅需几十纳秒，自带采样保持能力；
缺点：硬件开销极大，$n$位需要$2^n-1$个比较器与触发器，位数提升后硬件规模指数增长；转换精度由分压电阻、参考源精度决定。
4. \textbf{典型产品}
AD9012（8位TTL工艺）、AD9002（8位ECL工艺）、AD9020（10位TTL工艺）。

\subsubsection{反馈比较型ADC}
核心原理：将内部DAC输出的模拟电压和输入模拟电压对比，不断调整DAC输入数字量，直到二者逼近相等，最终数字量即为转换结果。
1. \textbf{计数型ADC}
由比较器、DAC、计数器、脉冲源、门控逻辑、输出寄存器组成：
- 清零后计数器从0开始逐脉冲累加，DAC输出阶梯电压逐步抬升；
- 当DAC输出超过输入电压时比较器翻转，关闭时钟门控，计数器数值锁存输出；
缺点：最长转换时长为$(2^n-1)T_{CP}$，速度很慢，仅适合低速场景。

2. \textbf{逐次渐近型ADC}
借鉴天平称重思路，从最高位到最低位逐位试探、逐位锁定，$n$位ADC仅需$n$次比较即可完成转换：
- 先将寄存器最高位置1，后续位清零，DAC输出试探电压和输入对比；
- 若DAC输出小于输入则保留该位1，否则清零，依次试探所有位；
优点：硬件规模远小于并联型、速度远快于计数型，是工业最通用的ADC架构。

\subsection{间接型A/D转换器——双积分型ADC}
\subsubsection{电路组成}
包含积分器、电压比较器、计数器、时序逻辑控制、时钟脉冲源五大模块，属于V-T变换型ADC。
\subsubsection{两步积分工作流程}
1. \textbf{第一次定时采样积分}：固定时长$T_1$对输入模拟电压$v_I$积分，积分输出：
\[
v_O=-\frac{T_1}{RC}v_I
\]
2. \textbf{第二次反向基准积分}：切换对固定负基准电压$-V_\text{REF}$积分，积分器从负电压往零回落，记录回落至0的时长$T_2$，满足：
\[
\frac{T_1}{RC}v_I=\frac{T_2}{RC}V_\text{REF} \implies T_2=\frac{T_1}{V_\text{REF}}v_I
\]
最终计数结果$D=\frac{T_2}{T_C}=\frac{2^n}{V_\text{REF}}v_I$，数字量和输入电压严格成正比。
\subsubsection{优缺点}
优点：抗工频干扰能力极强、对RC参数和时钟频率稳定性要求低、温漂小、转换精度极高；
缺点：转换速度慢，每秒仅能完成几十次转换，适合工业低速精密测量场景。
\subsubsection{典型集成芯片}
ADC-EK8B（8位二进制）、ADC-EK10B（10位二进制）、MC14433（$3\frac{1}{2}$位BCD码）、CB7106/CB7107（数码管驱动专用）。

\subsection{V-F变换型ADC}
先把模拟电压转换成与之成正比的脉冲频率，再在固定时间窗口内对频率脉冲计数，计数结果即为对应数字量；抗干扰性能优秀，适合远距离遥测场景。

\subsection{A/D转换器性能指标}
\subsubsection{分辨率}
用输出二进制/十进制位数表征，$n$位二进制ADC可区分$2^n$级输入电压，能分辨的最小输入电压为满量程的$\frac{1}{2^n}$，位数越高分辨率越强。
\subsubsection{转换误差}
实际输出数字量和理论理想值的最大偏差，工程要求误差不超过$\pm\frac{1}{2}\text{LSB}$；提升分辨率不能直接提高精度，必须让转换误差小于量化误差上限。
\subsubsection{转换速度}
用单次转换时间衡量：并联比较型速度最快、逐次渐近型次之、双积分间接型最慢；选型需要综合系统速度、精度、输入范围、功耗等多维度需求。

\subsection{集成A/D转换器及其应用}
\subsubsection{ADC0804集成芯片详解}
1. \textbf{基础特性}
采用CMOS工艺制作的逐次渐近型8位ADC，转换时长$100\mu\text{s}$，输入模拟电压范围$0\sim5\mathrm{V}$，片内自带输出锁存寄存器，可直接对接微机CPU数据总线，无需额外接口逻辑。

2. \textbf{引脚功能定义}
\begin{itemize}
\item $\boldsymbol{CS}$：片选信号输入，低电平有效，有效时芯片被选中启动工作；
\item $\boldsymbol{RD}$：读信号输入，低电平有效，$\boldsymbol{CS}$与$\boldsymbol{RD}$同时为低可读取转换结果；
\item $\boldsymbol{WR}$：写信号输入，低电平有效，$\boldsymbol{CS}$与$\boldsymbol{WR}$同时为低启动模数转换；
\item $\boldsymbol{CLKR}$：内部时钟发生器外接电阻端；
\item $\boldsymbol{CLKIN}$：外部时钟脉冲输入端；
\item $\boldsymbol{INTR}$：转换结束中断输出，低电平有效，转换完成后自动拉低，可作为CPU中断请求信号；
\item $\boldsymbol{V_{IN+}、V_{IN-}}$：差分模拟电压输入端，可接入单极性或双极性模拟信号；
\item $\boldsymbol{AGND}$：模拟信号地；$\boldsymbol{DGND}$：数字信号地；
\item $\boldsymbol{V_{REF}/2}$：参考电压分压输入端；
\item $\boldsymbol{D_7\sim D_0}$：8位三态并行数字输出端，可直接挂载微机总线；
\item $\boldsymbol{V_{CC}}$：芯片电源供电端。
\end{itemize}

3. \textbf{工作时序规范}
\begin{itemize}
\item 启动转换：$\boldsymbol{CS}$与$\boldsymbol{WR}$同时拉低，上升沿后开始转换，$100\mu\text{s}$内完成模数转换，转换结束后$\boldsymbol{INTR}$自动变低；
\item 读取结果：$\boldsymbol{CS}$与$\boldsymbol{RD}$同时拉低，内部三态门打开，数据输出到总线；电平恢复高电平后三态门回到高阻态。
\end{itemize}

4. \textbf{使用关键注意事项}
\begin{itemize}
\item 接地设计：模拟地（AGND）与数字地（DGND）需要分别布线，仅在电源处单点共接，避免数字开关噪声干扰模拟前端；
\item 零点与满量程调节：通过配套引脚外接微调电阻，修正零点偏差和满量程输出偏差；
\item 参考电压调节：根据输入量程配置参考电压，保障转换线性度；
\item 总线时序配合：严格匹配微机读写时序，防止数据采样出错。
\end{itemize}

5. \textbf{典型应用架构}
可直接挂接微机系统控制总线，配合RAM实现数据缓存；模拟信号送入ADC0804完成转换后，数字量上传至微处理器控制器，实现多路采集、数据处理与外设联动控制。